% Teaching manuscript. Edit this TeX source; Markdown is compatibility navigation only.

\StudyChapter{CS5222-Tutorial01}{Tutorial 1 · Switching, circuits and sharing｜把链路模型算清楚}
\StudyTarget{CS5222-Tutorial01__tutorial-1-switching-circuits-and-sharing}


\StudyRef{CS5222-Chapter01}{switching}{Chapter 1} · \StudyRepoRef{CS5222/course-notes/README.md}{}{课程目录} · \StudyRef{CS5222-Tutorial02}{root}{下一份Tutorial}


这份练习比较两种共享链路的方式。先为包画出逐跳时间轴，再为电路数清占用的边，最后用同时活跃人数解释统计复用。每题先固定模型条件，再计算容量或时间。

按当前 原题PDF（3页） 的Q1–4逐题讲解，并对照 教师解答PPT（10页）。下文“教师结论”指该文件；推演、时间表和变式由整理者补充。所有选择题的原选项在原页图中，正文给出推理步骤。


\textbf{可选先修：} \StudyRef{Foundations-NetworkBasics}{units}{bit/byte与速率}、\StudyRef{Foundations-NetworkBasics}{timeline}{流水线时间轴}、\StudyRef{Foundations-MathForML}{combinations}{组合概率}。


\begingroup
\small\setlength{\tabcolsep}{3pt}
\renewcommand{\arraystretch}{1.22}
\begin{longtable}{@{}>{\raggedright\arraybackslash}p{\dimexpr0.171726\linewidth-1.50000\tabcolsep\relax}>{\raggedright\arraybackslash}p{\dimexpr0.171726\linewidth-1.50000\tabcolsep\relax}>{\raggedright\arraybackslash}p{\dimexpr0.222384\linewidth-1.50000\tabcolsep\relax}>{\raggedright\arraybackslash}p{\dimexpr0.434164\linewidth-1.50000\tabcolsep\relax}@{}}
\toprule
原题 & 优先级 & 难度/卡点 & 本题检验；依据 \\
\midrule\endfirsthead
原题 & 优先级 & 难度/卡点 & 本题检验；依据 \\
\midrule\endhead
\bottomrule\endfoot
Q1 & 核心必会 & 中；不同包可以重叠 & 画存储转发时间轴；Tutorial 1教师解答，第3–4张幻灯片 \\
Q2 & 核心必会 & 中；每次连接占用哪些边 & 给出可行路由与容量上界；Tutorial 1教师解答，第6–7张幻灯片 \\
Q3 & 核心必会 & 低；单电路速率 & 单位换算与建立时延；Tutorial 1教师解答，第9张幻灯片 \\
Q4 & 核心必会 & 中；独立与组合数 & 容量和人数分布；Tutorial 1教师解答，第10张幻灯片 \\
\end{longtable}
\endgroup


学习优先级依据当前原题与教师解答。






\StudyHeading{1}{Q1. Store-and-forward｜先等整包到，再转发}
\StudyTarget{CS5222-Tutorial01__q1}
\StudyTarget{CS5222-Tutorial01__q1-store-and-forward}


\textbf{完整条件 / Conditions:} 一个源主机、一个分组交换机、一个目的主机；两条链路速率为R\textsubscript{1}、R\textsubscript{2}，包长L bit。存储转发，忽略排队、传播和处理时延。b问改为相同速率R，P=3个L-bit包跨N=2条链路。 / One switch connects a sender and receiver through links R\textsubscript{1} and R\textsubscript{2}. Packets contain L bits. Use store-and-forward and ignore queueing, propagation and processing. In part b, three equal packets cross two equal-rate R links.

\StudyImage{CS5222/course-notes/assets/tutorial01-01.png}
\StudyCaption{Tutorial 1原题页：Q1按链路次序排列每个包的开始发送与接收完成事件。}



\textbf{a的推演：} t=0开始发送；t=L/R\textsubscript{1}时整包抵达交换机；交换机此时才开始第二段，又花L/R\textsubscript{2}。因此教师结论为
\[
d=L/R_1+L/R_2\quad\text{（选C）}.\tag{T1.1}
\]


不是 \(L/(R_1+R_2)\)，因为这两条链路是先后经过，不是并行拼成一条更快链路。


\textbf{b的推演：} 令T=L/R，把每段链路当成一次只能发送一包的工位：


\begingroup
\small\setlength{\tabcolsep}{3pt}
\renewcommand{\arraystretch}{1.22}
\begin{longtable}{@{}>{\raggedright\arraybackslash}p{\dimexpr0.333333\linewidth-1.33333\tabcolsep\relax}>{\raggedright\arraybackslash}p{\dimexpr0.333333\linewidth-1.33333\tabcolsep\relax}>{\raggedright\arraybackslash}p{\dimexpr0.333333\linewidth-1.33333\tabcolsep\relax}@{}}
\toprule
时段 & 源\ensuremath{\to}交换机 & 交换机\ensuremath{\to}目的 \\
\midrule\endfirsthead
时段 & 源\ensuremath{\to}交换机 & 交换机\ensuremath{\to}目的 \\
\midrule\endhead
\bottomrule\endfoot
0到T & 包1 & 等待整包1 \\
T到2T & 包2 & 包1 \\
2T到3T & 包3 & 包2 \\
3T到4T & 无 & 包3 \\
\end{longtable}
\endgroup


最后整包到达时间4T=4L/R，选D；与教师解答第4页 \((P+1)L/R\) 一致。一般等速N跳为 \((P+N-1)L/R\)，但必须保留本题全部简化条件。


\textbf{English answer:} (a) L/R\textsubscript{1}+L/R\textsubscript{2}. (b) 4L/R. Store-and-forward constrains each packet, but separate links can transmit different packets simultaneously.


\textbf{变式 / Transfer:} 5个包、3条同速链路，同样忽略其他时延，多久？ / Five packets cross three equal-rate links under the same assumptions. Find completion time.



\begin{StudySolution}{答案 / Answer}


7L/R。第一包3L/R，之后4个每隔L/R到。 / 7L/R: three intervals for the first packet, then four further intervals.


\end{StudySolution}






\StudyHeading{1}{Q2. Circuit capacity｜数连接，也要数它消耗了几条边}
\StudyTarget{CS5222-Tutorial01__q2}
\StudyTarget{CS5222-Tutorial01__q2-circuit-capacity}


\textbf{条件 / Conditions:} A、B、C、D顺时针围成四边环，每条边4条电路，一次连接在路径每条边占一条电路。 / Four switches A,B,C,D form a clockwise ring; each link supports four circuits. A connection uses one circuit on every traversed link.

\StudyImage{CS5222/course-notes/assets/tutorial01-02.png}
\StudyCaption{Tutorial 1原题页：Q2一次连接占用其路径上的每一条边，容量不能只按端点数。}



\textbf{a 最大任意连接数 / Maximum arbitrary simultaneous connections:} 让每条边只承载相邻两点之间的连接，每边4个、4条边，总16，选C。上界也为16，因为全网总共16份“边电路”且每连接至少占一份，所以这个构造达到上界。


\textbf{b 全部A到C / All connections from A to C:} 每个至少两跳，可分别走A–B–C与A–D–C，各4，总8，选B。A出发仅8份电路，给出同样的上界；不是16\ensuremath{\div}4，因为每连接没有绕完整一圈。


\textbf{c 四个A–C加四个B–D / Four A–C plus four B–D:} 可行，选B（Yes）。按教师方案：A–C中2个经B、2个经D；B–D中2个经A、2个经C。


\begingroup
\small\setlength{\tabcolsep}{3pt}
\renewcommand{\arraystretch}{1.22}
\begin{longtable}{@{}>{\raggedright\arraybackslash}p{\dimexpr0.333333\linewidth-1.33333\tabcolsep\relax}>{\raggedright\arraybackslash}p{\dimexpr0.333333\linewidth-1.33333\tabcolsep\relax}>{\raggedright\arraybackslash}p{\dimexpr0.333333\linewidth-1.33333\tabcolsep\relax}@{}}
\toprule
路径 & 连接数 & 占用的边 \\
\midrule\endfirsthead
路径 & 连接数 & 占用的边 \\
\midrule\endhead
\bottomrule\endfoot
A–B–C & 2 & AB、BC \\
A–D–C & 2 & AD、DC \\
B–A–D & 2 & BA、AD \\
B–C–D & 2 & BC、CD \\
\end{longtable}
\endgroup


逐边相加AB=4、BC=4、CD=4、DA=4，不超过容量。不要只说“总共8小于16，所以可以”，那忽略了路径可能挤在同一条边。


\textbf{English answer:} (a)16; (b)8; (c)Yes, split each opposite-pair demand evenly between its two two-hop paths. Every link then carries exactly four circuits.


\textbf{变式 / Transfer:} 每边只有3条电路，全部A–C最多多少？ / With three circuits per link, what is the maximum for A–C connections only?



\begin{StudySolution}{答案 / Answer}


6，两条边不重叠路径各3。 / Six, with three on each disjoint two-hop path.


\end{StudySolution}






\StudyHeading{1}{Q3. FDM and setup｜总速率不等于你的那一份}
\StudyTarget{CS5222-Tutorial01__q3}
\StudyTarget{CS5222-Tutorial01__q3-fdm-and-setup}


\textbf{条件 / Conditions:} 文件160000 bit；电路网络用12个等宽频带，总速率1.536 Mbps，建立端到端电路需600 ms，随后才发文件。求总时间。 / A 160000-bit file uses one of twelve equal-rate FDM circuits sharing 1.536 Mbps. Circuit setup takes 600 ms before transmission. Find total time.


每电路速率 \(1.536\times10^6/12=128000\)bit/s。发送用 \(160000/128000=1.25\)s，加setup0.6s，共\textbf{1.85s，选A}。教师解答第9页相同。题目仅计建立和发送时延。


\textbf{English answer:} Each circuit supplies 128 kbps; 1.25 s transmission plus 0.60 s setup gives 1.85 s.


常见错法：直接除总速率1.536Mbps；把600ms当600s；把160000bit再乘8（本来就已经是bit）。


\textbf{变式 / Transfer:} 同电路发送320000bit、setup不变，总时间？ / Double the file size to 320000 bits without changing setup or circuit rate.



\begin{StudySolution}{答案 / Answer}


2.5+0.6=3.1s。 / 3.1 s. 文件翻倍只让发送部分翻倍，不把建立时延翻倍。


\end{StudySolution}






\StudyHeading{1}{Q4. Statistical multiplexing｜算人数前，先说是否独立}
\StudyTarget{CS5222-Tutorial01__q4}
\StudyTarget{CS5222-Tutorial01__q4-statistical-multiplexing}


\textbf{条件 / Conditions:} 共享3Mbps；每人活跃时需150kbps，每人10\%时间活跃。a用电路交换；b/c用分组交换，c共有120用户。 / Users share 3 Mbps, each requiring 150 kbps when active, with activity fraction 10\%. Part a uses circuits; parts b/c use packet switching, with 120 users in c.

\StudyImage{CS5222/course-notes/assets/tutorial01-03.png}
\StudyCaption{Tutorial 1原题页：Q4先区分预留容量与同时活跃人数，再计算超额概率。}



a 每用户预留150kbps，可容纳 \(3\times10^6/(150\times10^3)=20\)人，选B。不把它再除活跃比例，电路资源已经预留。


b 随机时刻某用户活跃概率0.1，选A。把长期活跃比例作为该随机时刻概率，是题目的稳态抽样约定。


c 在\textbf{用户独立、每人同为p=0.1}的模型下，K\textasciitilde{}Binomial(120,0.1)，恰n人活跃概率
\[
P(K=n)=\binom{120}{n}(0.1)^n(0.9)^{120-n},\quad n=0,1,\ldots,120.\tag{T1.2}
\]


对应原选项B（其因子顺序先写0.9）。组合数选出哪些人在发；0.1乘n次、0.9乘其余120\ensuremath{-}n次。原题没有清楚写独立，教师二项解答隐含这个假设，不能遗漏。


\textbf{English answer:} (a)20 users; (b)0.1; (c)choose(120,n)·0.1ⁿ·0.9\textsuperscript{1}\textsuperscript{2}\textsuperscript{0}\textsuperscript{-}ⁿ, assuming independent user activity.


平均活跃人数Np=12并不意味着始终12人；容量20也不保证任何时刻都够用。补充：超额事件为K>20，应该求从21起的尾概率，不是只算K=20或K=21。超额需求可用同一个人数分布继续计算。


\textbf{变式 / Transfer:} 3个独立用户、p=0.1，容量可同时容纳2人，超额概率？ / Three independent users have p=0.1 each; capacity supports two. Find excess-demand probability.



\begin{StudySolution}{答案 / Answer}


三人全活跃，0.1\textsuperscript{3}=0.001。 / All three must be active, giving 0.001.


\end{StudySolution}



\StudyHeading{1}{回查与巩固}
\StudyTarget{CS5222-Tutorial01___1}


所有数值与环路逐边占用由 \StudyRepoRef{CS5222/course-notes/tools/network_examples.py}{}{network\_examples.py} 复算。原题PDF p1=Q1/Q2引入，p2=Q2/Q3引入，p3=Q3继续/Q4；Tutorial 1教师解答，第1–2张幻灯片背景、3–4Q1、5电路图、6–7Q2、8FDM/TDM、9Q3、10Q4均已对应。


复习使用\href{https://crazyshout.github.io/micro-course/cards.html\#CS5222-N071}{现有Tutorial1卡}，首选N071/N072（存储转发）、N076（FDM）、N077（人数与容量），以及\StudyRef{CS5222-Chapter01}{qa}{Chapter1}相关题。学习完应能用英文解释 \emph{why}，不只记住C/D/B这些选项字母。
